\documentclass[11pt]{article}

\usepackage[a4paper,margin=25mm]{geometry}
\usepackage[T1]{fontenc}
\usepackage[utf8]{inputenc}
\usepackage{lmodern}
\usepackage{microtype}
\usepackage{amsmath,amssymb,mathtools}
\usepackage{amsthm}
\usepackage{array}
\usepackage{booktabs}
\usepackage{enumitem}
\usepackage{tabularx}
\usepackage{xcolor}
\usepackage{hyperref}

\definecolor{artifactblue}{HTML}{1F4E79}
\definecolor{artifactgreen}{HTML}{1E6B45}
\definecolor{artifactred}{HTML}{8B2F2F}

\hypersetup{
  colorlinks=true,
  linkcolor=artifactblue,
  urlcolor=artifactblue,
  pdftitle={A Mathematical Reading of TargetKeygenUniformXofLeaf.ec},
  pdfsubject={HAETAE uniform-XOF leaf range and frame verification}
}

\setlist[itemize]{leftmargin=1.5em,itemsep=0.25em,topsep=0.35em}
\setlist[enumerate]{leftmargin=1.7em,itemsep=0.25em,topsep=0.35em}
\renewcommand{\arraystretch}{1.18}
\setlength{\emergencystretch}{3em}

\theoremstyle{definition}
\newtheorem{definition}{Definition}
\theoremstyle{plain}
\newtheorem{theorem}{Theorem}
\newtheorem{corollary}{Corollary}
\theoremstyle{remark}
\newtheorem*{remark}{Remark}

\newcommand{\code}[1]{\nolinkurl{#1}}
\newcommand{\Z}{\mathbb Z}
\newcommand{\Wsixtyfour}{\mathbb W_{64}}
\newcommand{\Wthirtytwo}{\mathbb W_{32}}
\newcommand{\uint}[1]{\left\langle #1\right\rangle}
\newcommand{\word}[2]{\operatorname{word}_{#1}(#2)}
\newcommand{\Frame}{\operatorname{Frame}}
\newcommand{\Prefix}{\operatorname{Prefix}}
\newcommand{\Consumer}{\mathcal C}
\newcommand{\Leaf}{\mathcal U}

\newenvironment{resultbox}
  {\begin{center}\begin{minipage}{0.94\linewidth}\color{artifactgreen}\hrule\vspace{0.6em}\color{black}}
  {\vspace{0.6em}\color{artifactgreen}\hrule\end{minipage}\end{center}}

\newenvironment{assumptionbox}
  {\begin{center}\begin{minipage}{0.94\linewidth}\color{artifactblue}\hrule\vspace{0.6em}\color{black}}
  {\vspace{0.6em}\color{artifactblue}\hrule\end{minipage}\end{center}}

\newenvironment{boundarybox}
  {\begin{center}\begin{minipage}{0.94\linewidth}\color{artifactred}\hrule\vspace{0.6em}\color{black}}
  {\vspace{0.6em}\color{artifactred}\hrule\end{minipage}\end{center}}

\title{Uniform Rejection-Sampling Leaves in HAETAE\\
\large A mathematical reading of \code{TargetKeygenUniformXofLeaf.ec}}
\author{HAETAE reference EasyCrypt artifact}
\date{Artifact state checked 2026-07-23}

\begin{document}
\maketitle

\begin{abstract}
The extracted key-generation code fills a polynomial with 256 residues by
reading 16-bit candidates from a SHAKE128 path and rejecting candidates at
least $64513$.  The current EasyCrypt files prove the exact finite SHAKE128
stream-to-decoder relation through both generated leaves, coefficient range,
and confinement of writes to the assigned polynomial.  The bounded consumers
terminate unconditionally; each actual whole leaf terminates with probability
one under an explicit deterministic finite-prefix progress certificate and
exact seed/base/capacity/slice premises.  The development does not prove such
a certificate for every input or derive uniformity or independence from the
deterministic stream model.
\end{abstract}

\begin{resultbox}
\textbf{The result in one sentence.}
Let a destination contain either 2048 or 8192 32-bit slots.  If slots
$b,\ldots,b+255$ fit in that destination and the corresponding extracted leaf
returns, then those 256 slots contain integers in
$\{0,\ldots,64512\}$, while every byte outside those slots is unchanged.  Under
the corresponding exact seed and memory premises plus
\code{uniform_progress_prefix}, either actual leaf terminates with probability
one.
\end{resultbox}

\section{The Mathematical Question}

The program has two versions of the same sampler.  They differ only in the
capacity of the destination array.  Both versions start from a seed buffer,
an offset into that buffer, and a nonce.  They obtain bytes through an
extracted SHAKE128 code path, read consecutive pairs as 16-bit integers, and
keep a candidate precisely when it is less than
\[
  q=64513.
\]
They stop after accepting
\[
  n=256
\]
candidates.

The verification question answered by the two authored EasyCrypt theories is
deliberately narrower than full sampler correctness:

\begin{quote}
Assuming that the requested 256-word destination region fits, what can be
said about the returned words and about the rest of the destination array?
\end{quote}

The answer is independent of the contents of the byte stream.  The acceptance
test itself enforces the coefficient range, and the output index itself
enforces the frame property.  Consequently, these proofs can treat the
SHAKE-related calls without any semantic specification.

\section{Memory Model and the Two Capacities}

EasyCrypt models the relevant memories as fixed byte arrays.  A 32-bit word
occupies four consecutive bytes.  For an array $A$, write $\word{A}{i}$ for
the unsigned 32-bit word stored in slot $i$.

The names use two different units.  The suffix of an extracted leaf is its
capacity in \emph{32-bit slots}, whereas a type name such as
\code{BArray8192} gives its capacity in \emph{bytes}.

\begin{center}
\begin{tabular}{lccc}
\toprule
Leaf suffix & Destination type & Bytes & 32-bit slots \\
\midrule
\code{2048} & \code{BArray8192} & 8192 & 2048 \\
\code{8192} & \code{BArray32768} & 32768 & 8192 \\
\bottomrule
\end{tabular}
\end{center}

Thus ``the 2048 case'' below means an 8192-byte destination with 2048 word
slots, and ``the 8192 case'' means a 32768-byte destination with 8192 word
slots.  Correspondingly, specification identifiers ending in \code{8192}
refer to the small byte array, and identifiers ending in \code{32768} refer
to the large byte array.  A base is measured in word slots.  A seed offset is
measured in bytes.

Machine counters and bases belong to
\[
  \Wsixtyfour=\Z/2^{64}\Z.
\]
For $x\in\Wsixtyfour$, write $\uint{x}\in\{0,\ldots,2^{64}-1\}$ for
its unsigned representative.  Machine addition is modulo $2^{64}$; later we
will identify the condition that makes it agree with ordinary integer
addition.

\section{The Two Mathematical Predicates}

Fix a word capacity $m\in\{2048,8192\}$, so that the corresponding byte
array has length $4m$.  Let $A_0$ be a reference array, $A$ a current array,
and $b\in\Z$ an integer base.

\begin{definition}[The polynomial byte region]
The 256-word region starting at $b$ is the half-open byte interval
\[
  I_b=[\,4b,\ 4(b+n)\,).
\]
It has length $4n=1024$ bytes.
\end{definition}

\begin{definition}[Frame predicate]
The predicate
\[
  \Frame_m(A_0,A;b)
\]
holds when every valid byte outside $I_b$ has its original value:
\[
  \forall j\in\Z,\quad
  0\le j<4m\ \land\ j\notin I_b
  \ \Longrightarrow\ A[j]=A_0[j].
\]
The predicate intentionally says nothing about bytes inside $I_b$.
\end{definition}

\begin{definition}[Bounded accepted prefix]
For an integer $c$, define
\[
  \Prefix_m(A;b,c)
  \quad\Longleftrightarrow\quad
  \forall i\in\Z,\quad
  0\le i<c\ \Longrightarrow\
  0\le\word{A}{b+i}<q.
\]
The lower bound is automatic because \(\word{A}{b+i}\) is interpreted as an
unsigned word.  The EasyCrypt source therefore writes only the upper bound.
\end{definition}

The specification theory proves the elementary algebra needed later:
\begin{itemize}
\item $\Frame_m(A,A;b)$, and frame relations compose transitively;
\item writing one word at $b+s$, where $0\le s<n$, preserves the frame;
\item $\Prefix_m(A;b,0)$ holds vacuously; and
\item if $\Prefix_m(A;b,c)$, $0\le c<n$, and $0\le w<q$, then writing
  $w$ at $b+c$ establishes $\Prefix_m(A';b,c+1)$.
\end{itemize}

There is one necessary bridge between machine indices and integer indices.
For machine words $\beta,\gamma\in\Wsixtyfour$, if
\[
  \uint{\beta}+n\le m
  \qquad\text{and}\qquad
  \uint{\gamma}<n,
\]
then
\[
  \uint{\beta+\gamma}=\uint{\beta}+\uint{\gamma}.
  \tag{1}\label{eq:no-wrap}
\]
The capacity premise therefore prevents the machine computation of
\code{base + ctr} from wrapping modulo $2^{64}$ in any accepted-write branch.

\section{The Finite Consumer in Mathematical Form}

Let
\[
  \Consumer_m(A,b,c,B,L)=(A',c')
\]
denote the finite-buffer consumer for word capacity $m$.  Here $A$ is the
destination, $b,c\in\Wsixtyfour$, $B$ is a 1024-byte working buffer, and
$L\in\Wsixtyfour$ is the reported number of available bytes.

The extracted loop has the following normal form.  Its variables
$p$ (position) and $\ell$ (live flag) start at $0$ and $1$.

\begin{enumerate}
\item If $256\le_{\mathrm u}c$, set $\ell=0$.
\item Otherwise compute $r=L-p$ in $\Wsixtyfour$.  If
  $r<_{\mathrm u}2$, set $\ell=0$.
\item Otherwise read two bytes and form the little-endian candidate
  \[
    t=B[p]+2^8B[p+1]\in\{0,\ldots,2^{16}-1\},
  \]
  then replace $p$ by $p+2$.
\item If $t<q$, store the 32-bit representation of $t$ at slot $b+c$ and
  replace $c$ by $c+1$.  If $t\ge q$, leave $A$ and $c$ unchanged.
\item Repeat while $\ell\ne0$.
\end{enumerate}

This finite consumer does not promise to fill the polynomial.  It may return
with $c'<256$ because its current buffer has no complete two-byte candidate
left.  The whole leaf is responsible for obtaining more bytes and calling the
consumer again.

The separate lemmas \code{uniform_consume2048_ll} and
\code{uniform_consume8192_ll} prove that these finite consumers terminate on
every input.  This bounded-helper totality does not imply that the enclosing
rejection loop eventually collects 256 accepted candidates.

The two consumers have the same mathematical form.  The actual extracted
procedures are \code{__kp_poly_uniform_consume_2048} for $m=2048$ and
\code{__poly_uniform_consume} for $m=8192$.

\section{The Exact Partial-Correctness Theorems}

A Hoare statement
\[
  \{P\}\ f\ \{Q\}
\]
is read in this document as follows: if a call begins in a state satisfying
$P$ and the call returns, then its result satisfies $Q$.  It does not, by
itself, assert that the call returns.  This qualification applies to every
theorem in this section.

The next three results hold once for $m=2048$ and once for $m=8192$.

\begin{theorem}[Finite-consumer counter bound]
\label{thm:counter}
Let $c_0\in\Wsixtyfour$.  If the input counter is $c_0$ and
$\uint{c_0}\le n$, then, on return,
\[
  \uint{c_0}\le\uint{c'}\le n.
\]
No destination-capacity premise is needed for this counter-only statement.
\end{theorem}

\begin{theorem}[Finite-consumer frame preservation]
\label{thm:consumer-frame}
Suppose
\begin{align*}
  &\Frame_m(A_0,A;b_0),
  &&\uint{b}=b_0,\\
  &b_0+n\le m,
  &&0\le\uint{c}\le n.
\end{align*}
If
$\Consumer_m(A,b,c,B,L)=(A',c')$ returns, then
\[
  \Frame_m(A_0,A';b_0)
  \qquad\text{and}\qquad
  0\le\uint{c'}\le n.
\]
The input need not equal $A_0$; it is enough that earlier computation has
already preserved the same frame.
\end{theorem}

\begin{theorem}[Finite-consumer prefix extension]
\label{thm:consumer-range}
Suppose
\begin{align*}
  &\uint{b}=b_0,
  &&b_0+n\le m,\\
  &\Prefix_m(A;b_0,\uint{c}),
  &&0\le\uint{c}\le n.
\end{align*}
If
$\Consumer_m(A,b,c,B,L)=(A',c')$ returns, then
\[
  0\le\uint{c'}\le n
  \qquad\text{and}\qquad
  \Prefix_m(A';b_0,\uint{c'}).
\]
\end{theorem}

Now let
\[
  \Leaf_m(A,b,S,o,\nu)=A'
\]
denote the complete extracted leaf.  The seed buffer $S$ has 128 bytes, and
$b,o,\nu\in\Wsixtyfour$ are respectively the output base, seed offset, and
nonce.  Theorems~\ref{thm:leaf-range} and~\ref{thm:leaf-frame} are again
proved separately for both capacities.

\begin{theorem}[Whole-leaf range]
\label{thm:leaf-range}
If
\[
  \uint{b}=b_0
  \qquad\text{and}\qquad
  b_0+n\le m,
\]
then, whenever $\Leaf_m(A,b,S,o,\nu)$ returns an array $A'$,
\[
  \Prefix_m(A';b_0,n).
\]
Equivalently, all 256 returned target words lie in
$\{0,\ldots,64512\}$.
\end{theorem}

\begin{theorem}[Whole-leaf frame]
\label{thm:leaf-frame}
Under the same base and capacity premises, whenever
$\Leaf_m(A_0,b,S,o,\nu)$ returns an array $A'$,
\[
  \Frame_m(A_0,A';b_0).
\]
\end{theorem}

The range and frame properties appear as separate EasyCrypt lemmas.  Since
they concern the same deterministic procedures under the same capacity
premise, they may be read together as the following mathematical corollary.

\begin{corollary}[Returned polynomial block]
\label{cor:combined}
If the target region fits and the leaf returns, then
\[
  \forall i\in\{0,\ldots,255\},\qquad
  0\le\word{A'}{b_0+i}<64513,
\]
and
\[
  \forall j\in[0,4m)\setminus I_{b_0},\qquad A'[j]=A_0[j].
\]
\end{corollary}

\section{Why the Proof Is an Induction}

The EasyCrypt tactics implement ordinary invariant arguments.  Three
invariants contain essentially the entire mathematical proof.

\subsection{The counter invariant}

For a consumer that starts at $c_0$, use
\[
  \uint{c_0}\le\uint{c}\le256.
\]
A stop or rejection branch leaves $c$ unchanged.  An acceptance branch is
reached only while $c<_{\mathrm u}256$, and it changes $c$ to $c+1$.
Therefore it cannot pass 256, and the small value of $c$ prevents word
overflow.  This proves Theorem~\ref{thm:counter}.

\subsection{The frame invariant}

Use
\[
  \Frame_m(A_0,A;b_0)
  \quad\land\quad
  \uint{b}=b_0
  \quad\land\quad
  b_0+256\le m
  \quad\land\quad
  0\le\uint{c}\le256.
\]
Only an acceptance branch writes to the destination.  Its guard gives
$\uint{c}<256$, and equation~\eqref{eq:no-wrap} turns the machine index
$b+c$ into the integer slot $b_0+\uint{c}$.  That slot lies in the designated
256-word block, so the one-word frame lemma preserves the invariant.  This
proves Theorem~\ref{thm:consumer-frame}.

\subsection{The prefix invariant}

Use
\[
  \Prefix_m(A;b_0,\uint{c})
  \quad\land\quad
  \uint{b}=b_0
  \quad\land\quad
  b_0+256\le m
  \quad\land\quad
  0\le\uint{c}\le256.
\]
If a candidate is rejected, the prefix does not change.  If it is accepted,
the branch condition supplies $0\le t<q$.  The program writes $t$ into the
next slot and increments $c$, so the one-word prefix-extension lemma proves
the invariant for the new state.  This proves
Theorem~\ref{thm:consumer-range}.

\subsection{Lifting the invariant to the whole leaf}

The whole leaf initializes $c=0$, for which the prefix property is vacuous.
Every call to the finite consumer preserves it.  If the outer loop returns,
its false guard says $\uint{c}\ge256$, while the invariant says
$\uint{c}\le256$.  Thus $\uint{c}=256$, which changes a valid prefix into the
full conclusion of Theorem~\ref{thm:leaf-range}.  The frame argument lifts in
the same way, starting from the reflexive fact
$\Frame_m(A_0,A_0;b_0)$.

These range and frame proofs do not use a decreasing variant for the rejection
loop.  They derive their postconditions from the exit guard \emph{if exit
occurs}.  The finite consumers are unconditionally lossless by separate
bounded-variant proofs.  Separate pHoare theorems supply whole-leaf
losslessness under an explicit progress certificate; that stronger premise is
not needed for these functional postconditions.

\section{What the Whole Leaf Does Operationally}

The extracted target contains the following code path:
\begin{enumerate}
\item initialize a 200-byte Keccak state from 32 bytes beginning at the seed
  offset and two little-endian bytes of the nonce;
\item call the extracted SHAKE128 squeeze routine four times, obtaining
  $4\cdot168=672$ bytes;
\item call the finite consumer with counter zero;
\item while fewer than 256 values have been accepted, retain a possible
  one-byte tail, squeeze another 168 bytes, and call the consumer again.
\end{enumerate}

The nonce input has type $\Wsixtyfour$, but the initializer's two-iteration
loop absorbs only its low 16 bits.  The range and frame theorems impose no
premise on the seed offset or nonce.

The current stream theorems go beyond this operational reading.  Under an
in-bounds 32-byte seed slice, they relate the four-block prelude and every
later 168-byte squeeze to one exact finite SHAKE128 byte sequence, and relate
the stored prefix to the strict-$<64513$ candidates accepted from that
sequence.  Range and frame correctness remain valid independently of that
stronger stream interpretation.

\section{Certificate-Conditioned Whole-Leaf Termination}

For a fixed seed array $S$, seed offset $o$, nonce $\nu$, and block limit $L$,
let $C(r)$ be the number of strict-$<64513$ little-endian candidates among the
first $r$ consecutive 168-byte SHAKE128 blocks from the exact padded state for
$S[o:o+32]\Vert\operatorname{LE}_2(\nu)$.  The predicate
\code{uniform_progress_prefix S o nonce L}, written
$P_{\mathrm u}(S,o,\nu,L)$ below, is a deterministic certificate with two
parts:
\begin{enumerate}
\item $L\ge4$ and $C(L)\ge256$; and
\item for every $4\le r<L$, if $C(r)<256$, then $C(r)<C(r+1)$.
\end{enumerate}
It is a property of one concrete finite SHAKE128 byte sequence, not a
randomness assumption.  The strict-progress clause rules out a next block with
no accepted candidates at every certified incomplete full-block prefix.

\begin{theorem}[Actual uniform-leaf probability-one termination]
Fix an initial seed array $S$, seed offset $o$, nonce $\nu$, word base $b$,
and limit $L$.  Suppose the generated leaf receives exactly those seed,
offset, nonce, and base values, and
\[
  b+256\le m,
  \qquad \uint{o}+32\le128,
  \qquad P_{\mathrm u}(S,o,\nu,L),
\]
where $m=2048$ for the 2048-word-destination leaf and $m=8192$ for the
8192-word-destination leaf.  Then the corresponding actual leaf terminates
with probability one.
\end{theorem}

The two pHoare theorems are \code{uniform2048_leaf_progress_ll} and
\code{uniform8192_leaf_progress_ll}.  Because the generated leaves are
deterministic, probability one here is a losslessness statement under the
displayed precondition.  The development does not prove
\code{uniform_progress_prefix} for every seed, offset, and nonce, and it does
not lift per-leaf certificates through a parameterized caller in this file.
A separate parent refinement now lifts six certificates through the fixed
\(2\times3\) matrix caller and two through the fixed vector caller, then uses
them in a proof-only mode-2 sampler prefix.  It is not the full mode-2 parent.

\section{The Conditional Distribution Calculation}

\begin{assumptionbox}
\textbf{Conditional statement, not an EasyCrypt result in these files.}
Assume that the successive two-byte candidates are independent and uniformly
distributed in $\{0,\ldots,2^{16}-1\}$.  Under this extra assumption, the
accepted values are independent and uniform in $\{0,\ldots,q-1\}$, and the
rejection loop terminates almost surely.
\end{assumptionbox}

The calculation is elementary.  Exactly $q=64513$ of the $2^{16}=65536$
candidates pass the test, so
\[
  \Pr[t<q]=\frac{64513}{65536}\approx0.9843903.
\]
For every $a\in\{0,\ldots,q-1\}$,
\[
  \Pr[t=a\mid t<q]
  =\frac{1/65536}{q/65536}
  =\frac1q.
\]
The expected number of candidates needed for 256 acceptances would be
\[
  256\frac{65536}{64513}\approx260.06.
\]

This calculation explains why the implementation is shaped like a uniform
sampler.  The deterministic refinement now connects the extracted Keccak and
squeeze path to an exact SHAKE128 byte sequence, but it does not model that
sequence as random or prove the independence premise used in this
calculation.  The concrete progress certificate above supplies conditional
termination for one finite stream; it does not supply the probabilistic premise
or a universal certificate.

\section{Security Relevance and Proof Boundary}

The verified properties are useful components of a larger HAETAE proof:
\begin{itemize}
\item the 256 output words are canonical representatives modulo $q$;
\item accepted writes cannot cross the assigned polynomial boundary;
\item callers may compose the frame property when they allocate disjoint
  256-word regions; and
\item the same reasoning applies to both destination capacities.
\end{itemize}

They establish a well-formedness layer, not the sampler distribution required
by a cryptographic reduction:
\[
\begin{array}{c}
\text{extracted leaf}
\xrightarrow{\text{proved: range and frame}}
\text{well-formed polynomial memory}
\\[0.6em]
\text{well-formed polynomial memory}
\xrightarrow{\substack{\text{not proved here:}\\[-0.2em]
                        \text{XOF and distribution}}}
\text{ideal uniform vector in }(\Z/q\Z)^{256}.
\end{array}
\]

\begin{boundarybox}
\textbf{Precise boundary.}
The authored refinement file contains unary Hoare proofs about the target
program.  It does not contain an equivalence with an ideal random sampler.
The functional stream/range/frame claims described here are conditional on
return.
\end{boundarybox}

\noindent\begin{tabularx}{\textwidth}{>{\raggedright\arraybackslash}p{0.48\textwidth}X}
\toprule
Question & Status in these files \\
\midrule
Are all 256 returned target words below $q$? & Proved when the destination region fits. \\
Are bytes outside the target block unchanged? & Proved, under the same premise. \\
Is the finite consumer counter monotone and at most 256? & Proved, conditional on return. \\
Does the extracted Keccak-f transition and finite SHAKE128 path match the checked model? & Yes, through the exact finite-stream theorems. \\
Do the buffers form the intended finite XOF stream? & Yes, for the prelude and successive blocks covered by the leaf-stream contracts. \\
Do the actual whole leaves terminate under \code{uniform_progress_prefix} and the exact bounds? & Yes, with probability one. \\
Are coefficients uniform or independent? & Not proved. \\
Does rejection sampling terminate for every seed and nonce? & No universal progress certificate is proved. \\
Do the fixed mode-2 uniform callers terminate under bundled certificates? &
Yes, in the separate parent refinement.  The full mode-2 parent is not proved
lossless. \\
Is the leaf equivalent to an ideal sampler? & Not proved. \\
Are seed offsets in bounds and nonces unique? & Not proved. \\
Are physical memory accesses safe and destinations separated? & Not represented by these value-array Hoare theorems. \\
Is constant-time or speculative-execution security established? & Not proved by range and frame properties. \\
Is a HAETAE security-game advantage bounded? & Not addressed. \\
\bottomrule
\end{tabularx}

In particular, the capacity premise constrains the output base, not the seed
offset.  The theorems also do not prove that different callers use distinct
nonces or that the low-16-bit nonce encoding provides the paper's intended
domain separation.  Those are caller-composition and sampler-semantics
obligations.

\section{Translation Back to the EasyCrypt Files}

The three principal artifacts are:
\begin{center}
\footnotesize
\code{easycrypt/extract/keygen-sampler-callers/KeygenSamplerCallersTarget.ec}\\
\code{easycrypt/spec/KeygenUniformXofLeafSpec.ec}\\
\code{easycrypt/refinement/TargetKeygenUniformXofLeaf.ec}
\end{center}

The generated target contains the programs.  The authored specification
contains only mathematical predicates and their array lemmas.  The authored
refinement applies those predicates to the generated programs.

\begin{center}
\small
\begin{tabularx}{0.98\textwidth}{>{\raggedright\arraybackslash}p{0.37\textwidth}
                                  >{\raggedright\arraybackslash}p{0.33\textwidth}X}
\toprule
Mathematical object & EasyCrypt name & Source lines \\
\midrule
$q=64513$, $n=256$ & \code{uniform_q_i}, \code{uniform_poly_words_i} & 9--10 \\
Finite endpoint and progress certificate & \code{uniform_sufficient_prefix}, \code{uniform_progress_prefix} & 25--52 \\
Polynomial byte interval & \code{in_poly_bytes} & 67 \\
$\Frame_{2048}$, $\Frame_{8192}$ & \code{frame8192}, \code{frame32768} & 70--82 \\
$\Prefix_{2048}$, $\Prefix_{8192}$ & \code{bounded_prefix8192}, \code{bounded_prefix32768} & 84--92 \\
Frame algebra and one-word frame step & \code{frame*_refl}, \code{frame*_trans}, \code{frame*_set32} & 94--143 \\
Equation~\eqref{eq:no-wrap} & \code{base_counter_no_wrap*} & 150--169 \\
Machine-indexed frame step & \code{frame*_set32_word} & 170--207 \\
Prefix initialization and extension & \code{bounded_prefix*_zero}, \code{bounded_prefix*_succ}, \code{bounded_prefix*_set32_word} & 144--147; 208--291 \\
Certificate helper facts & \code{uniform_progress_prefix_*} & 462--510 \\
\bottomrule
\end{tabularx}
\end{center}

Here an asterisk denotes the separate 8192-byte and 32768-byte versions; it
is explanatory notation, not part of an EasyCrypt identifier.

\begin{center}
\small
\begin{tabularx}{0.98\textwidth}{>{\raggedright\arraybackslash}p{0.38\textwidth}
                                  >{\raggedright\arraybackslash}p{0.36\textwidth}X}
\toprule
Statement in this document & EasyCrypt lemmas & Refinement lines \\
\midrule
Theorem~\ref{thm:counter} & \code{consume2048_counter}, \code{consume8192_counter} & 218; 246 \\
Theorem~\ref{thm:consumer-frame} & \code{consume2048_frame}, \code{consume8192_frame} & 274; 324 \\
Theorem~\ref{thm:consumer-range} & \code{consume2048_range}, \code{consume8192_range} & 374; 429 \\
Unconditional finite-consumer termination & \code{uniform_consume2048_ll}, \code{uniform_consume8192_ll} & 1960; 2000 \\
Exact finite whole-leaf streams & \code{uniform2048_leaf_stream}, \code{uniform8192_leaf_stream} & 2040; 2704 \\
Theorem~\ref{thm:leaf-range} & \code{uniform2048_leaf_range}, \code{uniform8192_leaf_range} & 2687; 3351 \\
Theorem~\ref{thm:leaf-frame} & \code{uniform2048_leaf_frame}, \code{uniform8192_leaf_frame} & 3368; 3403 \\
Certificate-conditioned whole-leaf termination & \code{uniform2048_leaf_progress_ll}, \code{uniform8192_leaf_progress_ll} & 3438; 4047 \\
\bottomrule
\end{tabularx}
\end{center}

The generated-code correspondence is:
\begin{center}
\small
\begin{tabularx}{0.98\textwidth}{>{\raggedright\arraybackslash}p{0.55\textwidth}X}
\toprule
Generated procedure or code block & Target lines \\
\midrule
SHAKE128 squeeze of 168 bytes & 657--689 \\
8192-slot finite consumer, \code{__poly_uniform_consume} & 789--837 \\
Seed-and-nonce initializer, \code{__kp_shake128_init_seedbuf} & 838--891 \\
2048-slot finite consumer, \code{__kp_poly_uniform_consume_2048} & 946--991 \\
8192-slot whole leaf, \code{_kp_poly_uniform_at_seedbuf_8192} & 1079--1159 \\
2048-slot whole leaf, \code{_kp_poly_uniform_at_seedbuf_2048} & 1160--1242 \\
\bottomrule
\end{tabularx}
\end{center}

\section{Replaying the Machine-Checked Evidence}

From the \code{haetae-ref-easycrypt/} directory, run
\begin{center}
\code{./scripts/verify-keygen-sampler-callers-proof.sh}
\end{center}
This gate checks pinned sources and extraction drift, compiles the generated
target and the sampler specification/refinement manifest with
\code{-no-eco}, and scans the authored theories for proof holes and axiom
declarations.  It covers additional sampler-caller theories as well as the
uniform-XOF pair discussed here.

Build this document with two \LaTeX{} passes:
\begin{quote}
\begin{verbatim}
cd haetae-ref-easycrypt/latex
mkdir -p build
pdflatex -interaction=nonstopmode -halt-on-error \
  -output-directory=build TargetKeygenUniformXofLeaf.tex
pdflatex -interaction=nonstopmode -halt-on-error \
  -output-directory=build TargetKeygenUniformXofLeaf.tex
\end{verbatim}
\end{quote}
The PDF is written to
\code{build/TargetKeygenUniformXofLeaf.pdf}.

\section*{Conclusion}

The acceptance guard makes every stored candidate a canonical residue below
$q$, while the base-plus-counter index confines writes to one 256-word block.
The current proof also follows the exact finite SHAKE128 stream through both
whole leaves.  Their bounded consumers terminate unconditionally, and the
actual leaves terminate with probability one under their explicit deterministic
progress certificates.  Universal certificate existence, caller totality, and
the probabilistic uniformity/independence bridge remain separate obligations
for a complete HAETAE security argument; the fixed mode-2 callers and
proof-only first-attempt prefix now have certificate-conditioned totality.

\end{document}
